\chapter{Datasets \& Their Study}

\section*{Datasets as strategic policy objects}
\addcontentsline{toc}{section}{Datasets as strategic policy objects}

By unpacking the implications of literature exploring datasets as both constructed and strategic objects within ML value chains, I propose that establishing DIs to supply ML datasets would insert upstream, integrative policy surfaces into what are predominantly closed and opaque ML production pipelines. 

At present, AI governance focuses on harms and risks at the point of inference. Efforts have been made to delineate (in)appropriate deployment use cases for predictive and algorithmic decision-making (\cite{lum2016PredictServe, birhane2022AutomatingAmbiguityChallenges, aragon2022HumanCenteredDataScience}), to gatekeep access (\cite{shevlane2022StructuredAccessEmerging, gorwa2024ModeratingModelMarketplaces, shrestha2023BuildingOpensourceAI}), to mandate model requirements (\cite{bommasani2024FoundationModelTransparency}), and to implement algorithmic interventions altering system results towards unbiased outcomes (\cite{dwork2012FairnessAwareness, zhao2018GenderBiasCoreference, wang2022REVISEToolMeasuring, bellamy2018AIFairness360, googleGeminiVerge}). This inference-oriented approach deemphasises the role of inputs and their governance on AI outcomes. However, recognition that “there is no AI without data” (\cite{groger2021ThereNoAI, 2023WhatWeMean}) as a recent point of departure for the data-centric AI approach, broadens the scope of policy interventions beyond just system outputs, towards encompassing datasets. As AI progress continues to be driven by Deep Learning architectures learning the underlying structure of a dataset \textit{purely inductively}, Garbage In, Garbage Out (GIGO) - axiomatic of computer science since the 1960s - is now receiving due consideration with regards to data as representation for AI systems.

Data governance and AI governance represent interwoven realms - the former is prior and foundational to the latter (\cite{InterwovenRealmsData}). It follows then that the practice of dataset curation affords targeted intervention points into AI systems. Specifically, data enters AI systems through \textit{datasets}, and the data management literature’s focus on datasets over abstracted data, highlights their artificial, engineered and labour-intensive production processes. This insight suggests dataset curation's ability to motivate policy efforts around societal issues. Data availability constrains modelling effort towards tasks on which inference can be run (\cite{leavy2020DataPowerBias}) to the extent that supporting dataset development around a `task' or social issue supports technical efforts to address it. Further, that datasets are constructed, implies they can be strategically engineered towards socially-desirable outcomes (\cite{oala2023DMLRDatacentricMachine, verdegem2022DismantlingAICapitalism}). Considering this together with a framing of datasets as infrastructural to AI systems (themselves infrastructural to an `AI economy' (\cite{khan2022SubjectsStagesAI})) further implies that dataset governance ought to incorporate broad societal engagement and active public scrutiny, based on economic orthodoxy applied to public infrastructure projects elsewhere (\cite{singh2019DataDigitalIntelligence}). 

The data management literature employs a dual taxonomy for ML datasets, which I use to articulate divergent governance approaches. \textit{Training datasets} are used to teach models what they know - the majority are privately held, which serves as a barrier to entry for new market actors (\cite{spiekermann2021BigDataJustice, birkinbine2018CommonsPraxisCritical}), and an obstacle to standardising audit practices (\cite{raji2020ClosingAIAccountability, shrestha2023BuildingOpensourceAI}). Accordingly, governance priorities relate to distribution and access, as well as content and  harms contained therein. \textit{Benchmark datasets} are curated and stewarded by `task communities', and used to evaluate how well a model has learned to perform a technical `task' in a proxy-population setting (\cite{paullada2021DataItsDis}), they are mostly openly hosted on sites such as Kaggle and MLCommons, or integrated into development environments such as TensorFlow or PyTorch. \textcite{koch2021ReducedReusedRecycled} and \textcite{oala2023DMLRDatacentricMachine} highlight respectively the role of benchmark datasets in comparing between model types, and as estimation devices for how - and towards which ends - ML progress is being made. These task-oriented datasets connect policy challenges to fundamental ML research by providing the `right datasets for the right problems' - or don’t, leading to policy failure. \textcite{paullada2021DataItsDis} problematise how the development and task-focus of these strategic assets are driven by elite institutions, not always in alignment with the public interest. Governance priorities accordingly revolve around the incentives they structure within the ML research field.

The next generation of datasets tailored towards ML-specific purposes, open up new dimensions of quality across which DIs might create value. Efforts have been made to catalogue the technical requirements of ‘AI-ready’ datasets, characterised by dense metadata annotation (\cite{gebru2021DatasheetsDatasetsa, han2023NeglectedFreeLunch}), collaborative refinement, user preference and human feedback (\cite{kirstain2023PickaPicOpenDataset}), and evolution over time (\cite{schuhmann2022LAION5BOpenLargescale}). Research by \textcite{priestley2023SurveyDataQuality} into the burdensome process of developing safe, responsible ‘AI-ready’ datasets, and the extent to which those requirements are not currently met by profit-oriented ML developers, raise unanswered policy concerns that could inform the use and design of DIs as nodes in the ML value chain. Data in this developing framework is elevated from mere input resource - external to the model - to something like software code, iteratively evolving in conversation with model development (\cite{ jakubik2024DataCentricArtificialIntelligence}). By taking data work seriously and reconceptualising ML datasets as dynamic, labour-intensive artefacts, purposefully developed by stakeholders with agendas, reflective of and shaped by decisions made across their lifecycle, I hope to reveal previously unconsidered dimensions to DIs' value proposition for the ML community and wider society.

\section*{Datasets in AI governance}
\addcontentsline{toc}{section}{Datasets in AI governance}

\subsection*{Data as representation (and harms therein)}
\addcontentsline{toc}{subsection}{Data as representation (and harms therein)}

As an artefact of a decision to record a phenomena, data is biased to some irreducible degree. Nevertheless, as AI’s salience as a public issue develops, there is growing consensus around the need to consider its ramifications and pursue greater mitigation efforts. The academic study of bias (here operationalised as \textit{disparate results returned by sociotechnical systems across protected classes}) has spawned a robust quantitative fairness literature. Applied to the output of AI systems, this ranges from discussions of the use of ML in `fair affirmative action' (\cite{dwork2012FairnessAwareness}), to gender bias in commercial deployment of classification systems (\cite{buolamwini2018GenderShadesIntersectional}), and whose values are considered ground truths (\cite{crawford2021ExcavatingAIPolitics}). More recent \textit{data-centric} efforts have considered best practice in ML data collection (\cite{jo2020LessonsArchivesStrategies}), harms hidden inside training dataset content (\cite{birhane2021MultimodalDatasetsMisogyny}), bias creep in both label encodings (\cite{ghai2020MeasuringSocialBiases, dixon2018MeasuringMitigatingUnintended, miceli2020SubjectivityImpositionPower}) and outputs (\cite{abid2021PersistentAntiMuslimBias, gehman2020RealToxicityPromptsEvaluatingNeural, wolfe2022AmericanWhiteMultimodal}), as well as appropriate legal accountability frameworks (\cite{khan2022SubjectsStagesAI}). Similarly, various data-examination toolkits (\cite{bellamy2018AIFairness360, googleKnowYourData}), validation approaches (\cite{zhao2018GenderBiasCoreference}), and benchmarks (\cite{mazumder2023DataPerfBenchmarksDataCentric, zhao2018GenderBiasCoreference}) have been developed to evaluate datasets for biases. The salience of bias in AI is such that lively public debate around algorithmic decision-making in high-stakes public domains has sprung up, reaching into discussions of criminal justice and predictive policing (\cite{lum2016PredictServe}), hiring (\cite{hiringAlgo}), critical infrastructure and finance (\cite{lin2012NewInvestor}), and the ethics and (lack of scientific validity) in face-to-sexuality prediction (\cite{gehman2020RealToxicityPromptsEvaluatingNeural, wang2022REVISEToolMeasuring}), among other surveillance and inference-at-a-distance tasks. Nevertheless, it remains the case that \textit{in practice}, ML practitioners and researchers are drawn from CompSci-affiliated engineering domains, and that funding incentives (public and private) are driven by national strategic and business imperatives. These incentives motivate demonstrable system improvements over social science rigour - with deleterious implications for bias. DIs offer a commercial-and-policy-viable alternative. 

Most investigation of bias has thus far been directed at the system output level and the feedback loops engendered by algorithmic decision-making (see \textcite{abid2021PersistentAntiMuslimBias, buolamwini2018GenderShadesIntersectional, gehman2020RealToxicityPromptsEvaluatingNeural, wolfe2022AmericanWhiteMultimodal}). This is a convenient location at which to take measurements, but it has led to neglecting how these harms were encoded into the data inputs in the first place. Despite constraints (most training datasets are proprietary, and their size make them intractable to thoroughly scrutinise (\cite{bender2021DangersStochasticParrots})), a developing dataset auditing discipline is increasingly considering bias, class imbalance and other harms within these under-interrogated objects. This branch of investigation represents the most understood dimension of dataset-derived harms.

Recent (curtailed) efforts to audit ML training datasets reveal a myriad of representational harms (see \textcite{koch2021ReducedReusedRecycled} for a comprehensive overview). Considering only the most commonly used training datasets \textcite{buolamwini2018GenderShadesIntersectional} and \textcite{raji2021FaceSurveyFacial} independently demonstrate under-representation of darker-skinned subjects in over 100 facial analysis datasets, while \textcite{wilson2019PredictiveInequityObject} show predictive inequality in object-recognition systems - specifically skin tones of pedestrians - as does \textcite{devries2019DoesObjectRecognitiona} for regionally-sourced objects. Similarly \textcite{crawford2021ExcavatingAIPolitics} uncover millions of offensive labels within the ImageNet dataset; \textcite{prabhu2020LargeImageDatasets} demonstrate non-consensual and pornographic imagery throughout the TinyImages dataset; \textcite{zhao2018GenderBiasCoreference} identify under-representation of female pronouns in the OntoNotes dataset; \textcite{kreutzer2022QualityGlanceAudit} find male labels for images of both outdoor athletes (irrespective of identifiable gender) and where a person is too small for identification. At a higher level of abstraction, Prabhu articulates an “abattoir effect” and \textcite{jo2020LessonsArchivesStrategies} critique a culture of overly-permissive harnessing of massive datasets, with laissez-faire attitudes around their development. Alternatively, \textcite{paullada2021DataItsDis} identify routine, implicit invocation of particular, nonneutral viewpoints in the design of tasks and labelling heuristics, and point to pervasive undervaluation of data work.  

In contrast to other data-centric disciplines, ML practitioners (in aggregate) have been critiqued for lacking intentionality over the data they ingest, which has been linked to downstream harms. As such, datasets present a site of strategic contestation ripe for policy intervention. Operating at the point of collection, curation, distribution (and some preprocessing tasks), DIs and community-generated ML datasets offer a means of including a plurality of viewpoints, transparency and quality assurance mechanisms into this under-populated policy area.

\subsection*{Benchmarks as direction-setting devices for ML research}
\addcontentsline{toc}{subsection}{Benchmarks as direction-setting devices for ML research}

Beyond their contents, \textcite{koch2021ReducedReusedRecycled} assert that benchmark datasets (and associated metrics) \textit{direct} scientific ML progress and act as coordination devices for researchers around shared problems. Direction-setting is datasets' second strategic dimension. When institutionalising benchmarks, task communities endorse datasets as meaningful abstractions of a problem domain, and researchers are encouraged to make model choices that maximise performance accordingly. Alignment between benchmark datasets and “real world” tasks is thus critical; how alignment is measured, communicated and responded to offers a foothold for AI governance efforts. 

Concerningly, \textcite{koch2021ReducedReusedRecycled} also highlight field-wise concentration as a reduced number of elite (corporate, government, and academic) institutions shape the ML research agenda. Against this dynamic, establishing dedicated community-supported DIs for the continual development and maintenance of divers benchmark datasets, would better incorporate the interests and inputs of data-generating communities and associated stakeholders related to the `task' (technical or social) at hand. DIs able to make data available for beneficiary-approved tasks (and to withhold where not) would collectively form a decentralised, composite layer of context-appropriate democratic oversight procedures. As units within a system of commons-governance, a network of DIs could legitimately guide the appropriate structuring of ML research incentives, according to democratic principles. 

\subsection*{Datasets as bottleneck for leverage}
\addcontentsline{toc}{subsection}{Datasets as bottleneck for leverage}

The third strategic dimension of datasets for targeted AI governance is the oldest form of leverage: scarcity. Empirical research into ML scaling laws, suggests that while data and compute are today’s principal constraints, of the two, data will remain (see \textcite{kaplan2020ScalingLawsNeural} for detailed analysis). As such, data access, sovereignty, and producer protection are the bedrock of sustainable AI (\cite{curry2022DataSpacesDesign}). Society ought to be attentive to the ways that dataset owners and curators can be held accountable. Asserting more collective forms of control over access to datasets presents a point of democratic leverage in an otherwise concentrated industry structure. 

\section*{The practice of ML dataset development}
\addcontentsline{toc}{section}{The practice of ML dataset development}

It is not just ML datasets' contents that make them strategic points of entry for AI governance. Literature investigating (mal)practice in the process of dataset development identifies multiple sites of harm generation across both original data collection-generation, and ‘passive’ data scraping and mining.

This line of research serves as a meta-oriented parallel field of study to dataset auditing, focused on the generative processes behind the artefacts themselves. Insights generated here suggest these processes are amenable to the kinds of oversight improvement that DI governance offers. These insights could further inform policy-makers designing and implementing DIs by providing an architecture of existing pain points towards which to direct efforts.

\subsection*{Associated harms}
\addcontentsline{toc}{subsection}{Associated harms}

\subsubsection*{Licences}
\addcontentsline{toc}{subsubsection}{Licences}

Most major ML training datasets contain unlicensed, publicly available, internet-scraped data considered ‘all rights reserved’ by default. Despite the inclusion of robots.txt files or media clauses, the systematic inclusion of non-consenting sites is well demonstrated (\cite{khan2022SubjectsStagesAI}). Currently it is simply infeasible for ambitious commercial ML projects to train only on content for which open licences exist (\cite{chan2023ReclaimingDigitalCommons}). Benchmark datasets too are often drawn from a range of different sources, each with a different configuration of copyrights and permissions, although the problem is less severe and solution more manageable. Legal advocates have defended the use of copyrighted material under fair use and transformation clauses (\cite{sag2019NewLegalLandscape}). This remains controversial. Datasets curated through \textit{active} data collection-generation bypass some of these liability risks, but have been proven to generate a range of other harms (see \textcite{khan2022SubjectsStagesAI}). The decision to obtain or abide by appropriate licensing regimes is an irreducible source of tension and in near-constant violation by ML data users. 

DIs able to license entrusted data return value to data producers and subjects. While the economics of data licensing writ large remain intractable, a diverse ecosystem of bottom-up DIs could open up new areas of sufficiently overlapping interests between producers (data-generating individuals and communities) and consumers (ML developers) to support greater data sharing in an informed and licence-compliant manner (\cite{delacroix2019BottomupDataTrusts}). 

DIs could further choose to assume responsibility for ensuring their datasets are licence-compliant for a stated range of ML use cases - benefiting both parties equally; producers regain downstream data control, consumers gain compliance assurance. Data-stewarding DIs could draw on recent efforts to develop a gradation of experimental AI-oriented open data licences (\cite{dumas2023AI2ImpACTLicense}), to better represent their constituencies with greater granularity and freedom of movement, opening up dimensions of public value creation previously precluded by blunt licensing instruments (an alternative view is that recent proliferation of `open and responsible' AI licensing regimes erodes longstanding open source norms).

\subsubsection*{Documentation, consent \& distribution}
\addcontentsline{toc}{subsubsection}{Documentation, consent \& distribution}

Deficient documentation regarding datasets' consented use cases, the purpose of the original collection, or the conditions under which it was collected, are perennial challenges to ML developers. Investigations into major ML datasets such as the \textit{Common Crawl}, \textit{Faces in the Wild}, or \textit{ImageNet} reveal poor documentation and repurposing of personal data for uses that were not anticipated when consent for collection was given (\cite{khan2022SubjectsStagesAI}). By way of an early example, \textcite{radin2017DigitalNativesHow} was one of the first to document subsequent transpositions of a retracted (PIDD) dataset, tracing the harms derived from reusing data collected to refine algorithms that had nothing to do with that original understanding. 

Concerningly, most commercial models are released without publicly sharing training datasets, and a recent audit of 1800 datasets highlights diminishing efforts among AI developers to document datasets used (\cite{longpre2023DataProvenanceInitiative}). This is indicative of the fact that the release and distribution of datasets is typically not seen as part of the ML pipeline, leading social science critics contend that the ML field has yet to develop the data management practices required for the sensitive human data it is increasingly deployed to run inference on (\cite{raji2020ClosingAIAccountability, roh2019SurveyDataCollection, harvey2021ResearchersGoneWild}). 

Beyond the research ethics, deficient documentation and compromised data traceability is detrimental to responsible downstream model development (\cite{longpre2023DataProvenanceInitiative, gudivada2017DataQualityConsiderationsa}). Poor data genealogy hinders both the evaluation of data points' impact (\textit{leverage}) on systems’ outputs, and the ability to subsequently retract that data. These unintended afterlives (see accounts of derivative LFW, MS-Celeb-1M, and DukeMTMC datasets in \textcite{koch2021ReducedReusedRecycled}) suggest the need for targeted policy interventions around managing consent and retraction points earlier in the data cycle. DIs institutionalise such mechanisms - as potential agents of enforcement and representative single-points-of-contact around which to develop a predefined set of responsibilities for supporting recourse, and in generating the documentation and traceability programmes necessary to facilitate post-hoc remediation.

\subsubsection*{Market concentration}
\addcontentsline{toc}{subsubsection}{Market concentration}

The nature of data (alongside compute and human resources) make ML development a concentrated and opaque industry. As an economic and informational resource, data's utility lies in aggregation, and its collection benefits from economies of scale and cross-market leverage (\cite{cilauroINCREASINGACCESSDATA, coyle2021PotentialSocialValuea}). Data-driven web-2.0 markets favour anti-Schumpeterian ``winner-took-all" rentier incumbents conducting massive data collection across multiple markets on privatised platform infrastructure. (for CPE analyses, \textcite{spiekermann2021BigDataJustice, couldry2020CostsConnectionHow, meadwayCREATINGDIGITALCOMMONS}; for the new antitrust perspective, \cite{KhanAmazon2017}). Data's informational-relational aspects (\cite{viljoen2021RelationalTheoryData}) make for privacy concerns (\cite{zuboff}) and IP infringement (\cite{NYTGoogleOpenAIYoutube}) among other sensitivities - all induce growing levels of secrecy. In ML, the result is concentration and opacity in commercial data sourcing. This undermines the public’s ability to exercise decision-making power at strategic moments along the value chain, and puts data users in potential legal jeopardy (\cite{bommasani2024FoundationModelTransparency, bommasani2022OpportunitiesRisksFoundation}). 

Unlike antitrust and regulatory action, DIs for ML datasets address market concentration by leveraging data as an ‘AI bottleneck’ (while also accounting for its relational properties - unlike antitrust regulation). This efficiently negates the need for post-hoc distributive redress for the symptoms of market failure. Similarly, despite the proliferation of AI ethics standards (see \textcite{jobin2019GlobalLandscapeAI}), ethical guidelines are no substitute for addressing the structural factors underlying the concentration of AI resources. Commons-based DIs in contrast change the dynamics of power to offer \textit{integrative} industry restructuring - benefiting and compensating the majority of stakeholders while reassigning value closer to the true social costs/benefits.

\subsubsection*{Undervalued data work}
\addcontentsline{toc}{subsubsection}{Undervalued data work}

The final dimension of harm in ML dataset development lies in the production process itself. Prevailing ML data practises abstract away the human labour involved in dataset production (\cite{paullada2021DataItsDis}) andd wilfully ignore the subjective non-neutral judgements and labelling heuristics (biases) humans necessarily bring to the encoding table (\cite{jarrahi2023PrinciplesDataCentricAI}). 

For labelled data, developers purchase data products from vendors (Scale AI, Surge AI, Snorkel.ai etc) or directly engage gig crowdworkers (via platforms such as Amazon Mechanicl Turk (AMT)). Although semi-supervised labelling and automated annotation tools are growing (as problematised by \textcite{jarrahi2023PrinciplesDataCentricAI}), surveys of ML development practice highlight that cheap abundant gig labour remains the norm (\cite{priestley2023SurveyDataQuality}). Beyond the propensity for labour abuses under the platformisation of economic precarity (see \textcite{miceli2020SubjectivityImpositionPower}), this has deleterious implications for data quality. Scholars have demonstrated how AMT’s hierarchy positions data work and annotation as menial, leading to low quality results. This includes mistranslated text and mislabeled linguistic content (\cite{kreutzer2022QualityGlanceAudit, miceli2020SubjectivityImpositionPower}), divergences in judgements across (and within) annotator pools, (\cite{dixon2018MeasuringMitigatingUnintended, kreutzer2022QualityGlanceAudit}), subjective bias (\cite{ghai2020MeasuringSocialBiases, hube2019UnderstandingMitigatingWorker}), annotation that mischaracterise the real-world tasks the dataset supposedly represents (\cite{dixon2018MeasuringMitigatingUnintended}), and inadequate semantic linkages (\cite{WTFKnowledgeGraph}). These harms are further compounded by under-considered ad hoc data acquisition, resulting in distributional shifts, class imbalances, and hard-to-replicate phenomena (\cite{jarrahi2023PrinciplesDataCentricAI}).

Quality data processing is undersupplied at present. Privatised ML value chains are ill-equipped to consider its full social value, and unincentivised to internalise the full spectrum of harms generated. Commons-based governance for infrastructural ML data, where communities could own and profit (financially or otherwise) from the data they generate, would incentivise data quality, and encode diverse values into the necessarily subjective process. DIs offer an institutional template for structuring these elements of the data commons.

\subsection*{Taxonomies of dataset development}
\addcontentsline{toc}{subsection}{Taxonomies of dataset development}

By disaggregating the overlapping, multi-directional processes that constitute dataset development, I hope to map the identified technical and policy requirements of ML datasets to DI functionality. While implementation and granular design aspects of DIs are out of scope for this research, a working taxonomy of dataset development is necessary to later synthesis nine high-level principles according to the stage under considerations in the policy recommendations section. 

Across the stages of dataset development, decisions are made as to what kinds of data to collect, engineer and label, as well as how to steward and distribute that data (see \textcite{jernite2022DataGovernanceAge}). At present, these decisions remain mostly inscrutable and presented as purely technical considerations. However, given the harms detailed above, and the potential for ‘data cascades’ there is growing consensus that these decisions are neither neutral nor trivial, and that each decision point represents an opportunity for inserting greater democratic control (\cite{sambasivan2021EveryoneWantsModel}). DIs ought to be considerate of these decision points and designed accordingly.

There have been various overlapping taxonomies of sequencing dataset curation and their integration into the ML development process. As early as 1996, Fayyad et al. proposed nine stages. More recently, in their practitioner survey, \textcite{roh2019SurveyDataCollection} articulate three phases: acquisition (discovery, augmentation, generation), labelling, and improvement (of the data itself, or the model built upon it). This informs Priestley et al's (2023) framework for \textit{functionally grouping} tasks relating to use case and design, collection, and validation and maintenance. \textcite{desilva2022ArtificialIntelligenceLife} propose a chain of dynamic data \textit{requirements} to ML pipelines, from problem formulation, collection, cleaning, and annotation, to model training, evaluation and deployment, and inference. Other frameworks exist. All are idealised unidirectional flows of information, whereas the reality is messy iteration (\cite{munappy2022DataManagementProduction}). I draw on Priestley et al's framework in the policy recommendations below.

Much of this field of research comes from MLOps-oriented practitioners focused on identifying and disseminating best technical practice. Against this, a subset of papers consider encoding public values into technical data management decisions (\cite{priestley2023SurveyDataQuality, leavy2020DataPowerBias}), or identify labour abuse and exploitation along the supply chain ({\cite{paullada2021DataItsDis}). Similarly, climate policy concepts, like the `Twin-Transition’ analyse the physicality of the data cycle vis-a-vis carbon generation at each stage, and other climate concerns (see \textcite{TacklingClimateChangewithML2019}). For DIs  the relevant framework depends on their mission statement.
